\answer{ex:07:1}{$V=\kappa p/\bigl(\kappa p+(1-\kappa)(1-p)\bigr)$: $0.059$, $0.20$, $0.50$; медиана равна $0$, а точка~\eqref{eq:expectile} меняется с $\kappa$ плавно.}
\answer{ex:07:2}{$p=1/3$, $x=0.75$, $\sigma_r=0.35$, $V(0.2)=0.083$.}
\answer{ex:07:3}{$V=\sqrt\kappa/(\sqrt\kappa+\sqrt{1-\kappa})$, от $0.25$ до $0.75$; разброс $\propto\sqrt\eta$; у капель $V=0.25+0.5\kappa$, распределения различаются на $0.05$ у крайних нейронов, нужно $\eta\lesssim0.01$.}
\answer{ex:07:4}{(а)~$J^*=2\sqrt{C\bar R}$. (б)~Переплата $\bigl(k^{1/2}+k^{-1/2}\bigr)/2-1$: $6\,\%$ при $k=2$ и $25\,\%$ при $k=4$ --- хватает точности «в пределах двух раз».}
\answer{ex:07:5}{Всплеск площади $R(e^{-1}-e^{-2})\approx0.23R$ --- как $2.3\,\text{с}$ рампы; если дофамин сообщает $V$, всплеска нет, уровень просто ступенькой вырастает в $e$ раз.}
\answer{ex:07:6}{Событие сдвигает вес $\propto A\tau_f/(\tau_e+\tau_f)$, уровень даёт ошибку $s\tau_f$: $9\,\text{с}\le\tau_f\le17\,\text{с}$.}
\answer{ex:07:7}{Исходно $\Delta P=0.80$, $M=0.90$. (а)~$\Delta P=0.74$ ($-8\,\%$), $M=0.43$ ($-52\,\%$); (б)~$\Delta P=0.40$ ($-50\,\%$), $M=0.80$ ($-11\,\%$): двойная диссоциация.}
\answer{ex:07:8}{$N_{\text{эфф}}\to2+\rho^2N\approx10^6$ попыток: в $10^4$ раз меньше, чем с одним проводом, но утечка в $1\,\%$ всё ещё стоит миллиона попыток.}
